\section{Structured graph families and constructor conventions}
In addition to cycles $C_n$, complete graphs $K_n$, and hypercubes $Q_d$,
the benchmarks consider three Cartesian families,
\[
 C_n\mathbin{\square} C_m,\qquad C_n\mathbin{\square} P_m,\qquad P_n\mathbin{\square} P_m,
\]
and four corona families,
\[
 C_n\circ C_m,\quad C_n\circ P_m,\quad P_n\circ C_m,\quad P_n\circ P_m.
\]
A Cartesian product has vertices $(u,a)$, with an edge when one coordinate
is unchanged and the other is an edge of its factor graph. The corona
$G\circ H$ consists of one copy of $G$, a separate copy of $H$ for each
$u\in V(G)$, and edges from $u$ to every vertex in its corresponding copy
\cite{frucht}. If $G,H$ have $n,m$ vertices, respectively, then
\[
|V(G\circ H)|=n(m+1),\qquad
|E(G\circ H)|=|E(G)|+n|E(H)|+nm.
\]
The current implementation uses ordinary coronas, not edge or neighborhood coronas.

The implemented $GP(n,k)$ has edges $u_iu_{i+1}$, $u_iv_i$, and $v_iv_{i+k}$,
with indices modulo $n$ and $1\leq k<n/2$. It has $2n$ vertices and $3n$ edges.
This must be distinguished from $GPG(n)$ in Huang et al.\ \cite{huang},
which consists of two length-$n$ cycles joined by an arbitrary matching.
When $\gcd(n,k)=1$, the inner cycle of $GP(n,k)$ is connected and the graph
belongs to that class; when $\gcd(n,k)>1$, the step-$k$ edges form several
cycles. Sweeping $k$ also does not enumerate all matchings of $GPG(n)$.
Thus a $GP(n,k)$ sweep does not verify the full Georges--Mauro conjecture
on $GPG(n)$.
